\documentclass[11pt]{article}
\usepackage[margin=1in]{geometry}
\usepackage{hyperref}
\usepackage{enumitem}
\setlist[itemize]{leftmargin=1.5em}
\begin{document}
\title{Machine Learning Notebook Table of Contents}
\author{}
\date{}
\maketitle
\section*{Overview}
This document lists the topics covered in each notebook under the Machine Learning directory.
\section*{MLALGO001\_Linear Regression.ipynb}
\begin{itemize}
\item Linear Regression (OLS) - Practical Guide
\item Simple Linear Regression
\item Correlation and Covariance
\item MLR - Matrix Formulation
\item Regularization (Ridge, Lasso, Elastic Net)
\item Linear Regression Implementation: Numpy
\item Linear Regression Implementation: List
\item Standard Errors of Estimates
\item Estimate of the sd of the error term : RSE
\item R2 Statistic
\item Adjusted R2
\item AIC \& BIC Calculation
\item Hypothesis Testing
\item F - test in regression analysis
\item Confidence Intervals \& Prediction Intervals
\item OSL Violations
\item Residual Diagnostics
\item residuals vs fitted
\item Q-Q plot ( Qunatile - Quantile)
\item Shapiro-wilk test for normality
\item Kolmogorov-Smirnov: Normality
\item Anderson-Darling Test : Normality
\item One-sample t-test for mean of residuals=0
\item Leverage Plot
\item Scale-Location Plot
\item Breusch–Pagan Test for Homoscedasticity
\item Durbin-Watson Test for Autocorrelation of Residuals
\item VIF - Test for Multicollinearity
\item Diagnostic Test Class - Code Snippet
\end{itemize}
\section*{MLALGO002\_Logistic Regression.ipynb}
\begin{itemize}
\item Model Formulation
\item Linking Functions
\item Parameter Estimation using MLE
\item Optimization Talks
\item Canonical Exponential Form
\item Implementation
\item Key-Logistic Assumptions
\item Odds, log -odds, odds ratio - Interpretations
\item Overall Model Significance - Likelihood Ratio Test
\item Wald Test - Single Coefficient Significance Test
\item Comparison
\end{itemize}
\section*{MLALGO003\_KNN Algorithm.ipynb}
\begin{itemize}
\item K - nearest neighbor Algorithm
\item Implementation
\end{itemize}
\section*{MLALGO004\_Naive Bayes.ipynb}
\begin{itemize}
\item Naive Bayes Algorithm
\item Gaussian Naive Bayes ( continuous features)
\item Bernoulli Naive Bayes ( Binary features)
\item Multinomial Naive Bayes ( text, counts)
\item Mixed DataType
\item custom implementation:
\end{itemize}
\section*{MLALGO005\_Support Vector Machines.ipynb}
\begin{itemize}
\item Algorithm \& Formulation
\item Objective \& Optimization
\item Kernel Trick
\item Simple Implementation
\item Soft Margin
\item Non-linearly separable data
\item rbf kernel
\item Kernel : Mapping
\end{itemize}
\section*{MLALGO006\_Decision Trees.ipynb}
\begin{itemize}
\item Decision Tree Algorithm
\item What should be the split ?
\item CART Algorithm
\item Encoding Categorical
\item Stopping Criteria ( Why Trees Overfit )
\item Merits of Decision Trees
\item Demerits of Decision Trees
\item Properties of Decision Trees
\item Pre-Pruning Decision Trees
\item Cost - Complexity Pruning ( Post-Pruning)
\item Custom Decision Tree
\end{itemize}
\section*{MLALGO007\_Random Forest.ipynb}
\begin{itemize}
\item Ensemble Methods
\item Bagging
\item Boosting
\item Stacking
\item Voting
\item Random Forest
\item Bias Variance Trade off
\item Randomness in Random Forest
\item Out of Bag Score (OOB Score)
\item Proximity Matrix
\end{itemize}
\section*{MLALGO008\_Adaboost Algorithm.ipynb}
\begin{itemize}
\item Adaboost Algorithm
\end{itemize}
\section*{MLALGO009\_Gradient Boosting.ipynb}
\textit{No markdown headings found.}
\section*{MLALGO010\_XGBoost Algorithm.ipynb}
\begin{itemize}
\item XGBoost
\end{itemize}
\section*{MLALGO011\_Catboost\_lgbm.ipynb}
\textit{No markdown headings found.}
\section*{MLALGO012\_Kmeans Clustering.ipynb}
\begin{itemize}
\item Optimized Implementation usng Numpy
\item K-Prototype
\item Simple implementation using : list
\end{itemize}
\section*{MLALGO013\_DBSCAN.ipynb}
\textit{No markdown headings found.}
\section*{MLALGO014\_Hierarchical Clustering.ipynb}
\textit{No markdown headings found.}
\section*{MLALGO015\_Recommendation System.ipynb}
\begin{itemize}
\item Fundamental Notions
\item User-Based Collaborative Filtering
\item Item - Based Collaborative Filtering
\item Content - Based Filtering
\item Matrix Factorization
\end{itemize}
\section*{MLALGO016\_GMM.ipynb}
\begin{itemize}
\item Gaussian Mixture Models-Expectation Maximisation Algorithm
\item EM Algorithm
\item Gaussian Mixture Model
\end{itemize}
\section*{MLALGO017\_Isolation Forest.ipynb}
\begin{itemize}
\item Isolation Forest
\item Steps
\item Mathematical Formulation
\item Why subsampling is required ?
\item Implementation
\end{itemize}
\section*{MLALGO018\_PCA.ipynb}
\textit{No markdown headings found.}
\section*{MLFUND001\_Bias Variance Tradeoff.ipynb}
\textit{No markdown headings found.}
\section*{MLFUND002\_ShapExplanation.ipynb}
\begin{itemize}
\item Explainable AI - SHAP
\item SHAP - Shapley Additive Explanations
\item The interpretation of the Shapley value for feature value j is: The value of the j-th feature contributed \$\textbackslash\{\}large\textbackslash\{\}phi\_j\$to the prediction of this particular instance compared to the average prediction for the dataset.
\item Waterfall Plots
\item Decision Plots
\item Simple dependence plot
\end{itemize}
\section*{MLFUND003\_Regression Evaluation Metrics.ipynb}
\begin{itemize}
\item Regression Evaluation Metrics
\end{itemize}
\section*{MLFUND004\_Classification Evaluation Metrics.ipynb}
\begin{itemize}
\item Confusion Matrix Measurements
\item AUC \& Gini
\item **Multi Class Confusion Matrix**
\item Micro and Macro Measurements
\item KS Statistic (Kolmogorov-Smirnov)
\end{itemize}
\section*{MLFUND005\_Model Calibration.ipynb}
\begin{itemize}
\item Model Calibration
\item Steps in Model Calibration
\item Brier Score Loss
\item Platt Scaling or Isotonic Regression
\item Isotonic Regression
\item 1. Why We Need Isotonic Regression
\item 2. What Isotonic Regression Does
\item Block Averaging (Key Idea)
\item Pool-Adjacent-Violators (PAV) Algorithm
\item Interpretation
\end{itemize}
\section*{MLFUND006\_Feature Selection Methods.ipynb}
\begin{itemize}
\item Filter Methods In Feature Selection
\item Mutual Information
\item ANOVA For Feature Selection
\item Variation Inflation Factor - VIF
\item Weight of Evidence
\item Chi-Square Test of Independence
\item Wrapper Methods in Feature Selection
\item Forward Selection
\item Backward Selection
\item Stepwise Selection
\item Recursive Feature Elimination
\item Embedded Methods in Feature Selection
\item Lasso - L1 Regularisation
\end{itemize}
\section*{MLFUND007\_Feature Importance.ipynb}
\begin{itemize}
\item Tree-based importance
\item Permutation importance
\item Linear model coefficients
\item Drop-column
\item LIME
\item Partial Dependency Plot
\item ICE - Individual Conditional Expectation Plot
\end{itemize}
\section*{MLFUND009\_Population Stability Index.ipynb}
\begin{itemize}
\item Population Stability Index
\end{itemize}
\section*{MLFUND010\_Information Theory Topics.ipynb}
\textit{No markdown headings found.}
\section*{MLFUND011\_Loss\_Functions.ipynb}
\begin{itemize}
\item **Metrics from L-P Norm**
\item Entropy
\item Cross Entropy Loss
\item Joint Entropy
\item Mutual Information
\item KL Kullback-Leibler Divergence
\item Huber Loss
\end{itemize}
\section*{MLFUND012\_Covariance Matrix \& Eigen Values.ipynb}
\begin{itemize}
\item Covariance Matrix Transformation
\item Covariance Matrix as a linear Transformation
\item Applying covariance matrix transformation n times on the same dataset ;
\item EigenVectors are the Principal Components
\item 1. What are Eigenvectors and Eigenvalues?
\item 2. Conditions for Eigenvectors
\item 3. What Does This Have to Do with Linear Transformations?
\item 4. Eigenvectors and Covariance Matrices
\item The Principal Components:
\item 5. Proof Outline
\item Step 1: Covariance Matrix and Linear Transformation
\item Step 2: Eigenvalue Decomposition
\item Step 3: Maximizing Variance
\item 6. Conclusion
\item 7. Summary
\item Vector Projections
\item 1. What is Vector Projection?
\item 2. Geometric Interpretation of Projection
\item 3. Vector Projections in Principal Component Analysis (PCA)
\item Step-by-Step Role of Projections in PCA:
\item Example:
\item 4. Why Projections are Key in PCA?
\item 5. Orthogonality of Principal Components
\item 6. Relation to Singular Value Decomposition (SVD)
\item 7. Summary of Vector Projections in PCA
\item Transformation with the Inverse of the Covariance Matrix
\item 1. **Transformation with the Inverse of the Covariance Matrix**
\item 2. **Intuitive Explanation of the Transformation**
\item **What happens when we apply the inverse covariance matrix?**
\item 3. **Mathematical Viewpoint:**
\item 4. **Effect on Multivariate Gaussian Distribution**
\item Original Distribution (before applying \textbackslash\{\}( \textbackslash\{\}Sigma\textasciicircum{}\{-1\} \textbackslash\{\})):
\item After Transformation (whitening or normalization):
\item 5. **Summary of the Transformation with the Inverse Covariance Matrix:**
\end{itemize}
\section*{MLFUND013\_Shapley Values.ipynb}
\begin{itemize}
\item Shapley Values
\item Shapley Axioms
\item Machine Learning Application
\item Consistent method for feature attribution
\item Variants of SHAP
\end{itemize}
\section*{MLFUND014\_Cluster Evauation Metrics\_.ipynb}
\begin{itemize}
\item Silhouette Score for Clustering Evaluation
\item Davies-Bouldin Score
\item Calinski-Harabasz (CH) Score for Clustering Evaluation
\item Key Concept
\item Between-Cluster Dispersion
\item Within-Cluster Dispersion
\item Calinski-Harabasz Index Equation
\item Intuition Behind the Formula
\item Calculation Steps
\end{itemize}
\section*{MLFUND015\_Similarity Metrics.ipynb}
\begin{itemize}
\item Euclidean
\item Cosine Similarity
\item Mahalanobis Distance
\item Jaccard Similarity
\item Hamming Distance
\item Manhattan Distance
\item Minkowski Distance
\item Chebyshev Distance
\item Bray-Curtis Dissimilarity
\item Correlation Methods : Pearson,Spearman, KendalTau
\item Gower Score
\end{itemize}
\section*{MLFUND016\_Cluster Analysis Methods.ipynb}
\textit{No markdown headings found.}
\section*{MLFUND017\_Customer Segmentation.ipynb}
\begin{itemize}
\item Customer Segmentation
\item Basic Information used in segmenation
\item Extra Information
\item Types of Customer Segmentation Models
\item Benefits of Customer Segmentation
\item Cross-Selling and upselling
\end{itemize}
\section*{MLFUND018\_Lime.ipynb}
\begin{itemize}
\item Local Interpretable Model Agnostic Explanation
\item Lime algorithm
\item Sub-Modular pick Algorithm
\item Advantages/Disadvantages
\end{itemize}
\section*{MLFUND019\_Data Analysis.ipynb}
\begin{itemize}
\item correlation droppings
\item feature analysis
\item woe
\item chi-square selection
\item Entropy
\end{itemize}
\section*{MLFUND020\_data\_manipulation.ipynb}
\begin{itemize}
\item Multi Grouping - aggregation/sorting
\item Fetching the nth item after groupby/orderby
\item Groupby orderby - assigning ranks
\item Using groupby and apply together
\item Pivot Table Usage
\item apply,map
\item Grouping row values into a list
\item melt function
\end{itemize}
\end{document}